\section{广义 KKT 条件}

\label{sec:7-4}

到这里为止，我们已经有了三块拼图：第 \ref{sec:7-1} 节给出锥与对偶锥，第 \ref{sec:7-2} 节给出抽象标准型，第 \ref{sec:7-3} 节给出强对偶成立的资格条件。KKT 理论所做的，就是把这三块拼图与第 \ref{sec:6-4} 节的次微分最优性条件合成一个统一系统。有限维里这常常被写成若干向量等式；无穷维里真正需要小心的是，乘子所属的空间、法锥的表达方式以及互补松弛究竟对哪一个锥而言成立。

\subsection{KKT 系统的定义}

\begin{definition}[锥规划的 KKT 系统]\label{def:kkt-cone-system}
考虑定义~\ref{def:abstract-cone-program} 的原问题
\[
\inf_{x\in X}\{f(x):Ax-b\in -K\}.
\]
称 $(\bar x,\bar y^\ast)\in X\times Y^\ast$ 满足该问题的\emph{KKT 系统}，若
\[
A\bar x-b\in -K,
\]
\[
\bar y^\ast\in K^\ast,
\]
\[
0\in \partial f(\bar x)+A^\ast \bar y^\ast,
\]
\[
\langle \bar y^\ast,b-A\bar x\rangle=0.
\]
四个条件分别称为原可行性、对偶可行性、驻点条件和互补松弛。
\end{definition}

\begin{proposition}[锥松弛与对偶锥的互补配对]\label{prop:cone-complementarity-dual-pair}
设 $K\subset Y$ 为闭凸锥，$s\in K$，$y^\ast\in Y^\ast$。则下列两条等价：
\[
y^\ast\in N_{-K}(-s);
\]
\[
y^\ast\in K^\ast,\qquad \langle y^\ast,s\rangle=0.
\]
等价地，对任意 $b\in Y$，
\[
y^\ast\in N_{b-K}(b-s)
\quad \Longleftrightarrow \quad
y^\ast\in K^\ast,\ \langle y^\ast,s\rangle=0.
\]
\end{proposition}

\begin{proof}
先证明关于 $-K$ 的表述。若 $y^\ast\in N_{-K}(-s)$，则对任意 $k\in K$，由于 $-k\in -K$，有
\[
\langle y^\ast,(-k)-(-s)\rangle\le 0,
\]
即
\[
\langle y^\ast,s-k\rangle\le 0.
\]
取 $k=0$ 得 $\langle y^\ast,s\rangle\le 0$；取 $k=2s$ 又得 $\langle y^\ast,-s\rangle\le 0$，故
\[
\langle y^\ast,s\rangle=0.
\]
再代回上式，便得
\[
\langle y^\ast,k\rangle\ge 0,\qquad \forall k\in K,
\]
即 $y^\ast\in K^\ast$。

反之，若 $y^\ast\in K^\ast$ 且 $\langle y^\ast,s\rangle=0$，则对任意 $k\in K$，
\[
\langle y^\ast,(-k)-(-s)\rangle
=
\langle y^\ast,s\rangle-\langle y^\ast,k\rangle
=
-\langle y^\ast,k\rangle\le 0,
\]
从而 $y^\ast\in N_{-K}(-s)$。

最后注意到平移不会改变法锥的方向结构：
\[
N_{b-K}(b-s)=N_{-K}(-s).
\]
于是第二个等价关系随即成立。
\end{proof}

\subsection{广义 KKT 定理}

\begin{theorem}[无穷维锥规划的 KKT 定理]\label{thm:infinite-dimensional-kkt}
设 $X,Y$ 为 Banach 空间，$A\colon X\to Y$ 为有界线性算子，$K\subset Y$ 为闭凸锥，$f\colon X\to \mathbb R\cup\{+\infty\}$ 为 proper、凸且下半连续的泛函。假设第 \ref{sec:7-3} 节定义~\ref{def:generalized-slater-condition} 成立，并设 $\bar x\in X$ 与 $\bar y^\ast\in Y^\ast$ 分别为原问题和对偶问题的可行点。则下列两条等价：
\begin{enumerate}
  \item $\bar x$ 与 $\bar y^\ast$ 分别是原问题与对偶问题的最优解；
  \item $(\bar x,\bar y^\ast)$ 满足定义~\ref{def:kkt-cone-system} 的 KKT 条件。
\end{enumerate}
\end{theorem}

\begin{proof}
把原问题写成
\[
\inf_{x\in X}\{f(x)+\iota_{b-K}(Ax)\}.
\]
由第 \ref{sec:7-3} 节定理~\ref{thm:slater-strong-duality-cone}，该问题满足强对偶。于是可应用第 \ref{sec:6-4} 节推论~\ref{cor:fenchel-kkt-subdifferential}：$(\bar x,\bar y^\ast)$ 为原--对偶最优对，当且仅当
\[
-A^\ast \bar y^\ast\in \partial f(\bar x),
\qquad
\bar y^\ast\in \partial \iota_{b-K}(A\bar x).
\]
由第 \ref{sec:5-4} 节推论~\ref{cor:normal-cone-indicator}，
\[
\partial \iota_{b-K}(A\bar x)=N_{b-K}(A\bar x).
\]
再令
\[
s:=b-A\bar x.
\]
则原可行性就是 $s\in K$，而命题~\ref{prop:cone-complementarity-dual-pair} 给出
\[
\bar y^\ast\in N_{b-K}(A\bar x)
\quad \Longleftrightarrow \quad
\bar y^\ast\in K^\ast,\ \langle \bar y^\ast,b-A\bar x\rangle=0.
\]
因此，第 \ref{sec:6-4} 节的 Fenchel 次微分最优性条件恰好等价于
\[
A\bar x-b\in -K,\qquad
\bar y^\ast\in K^\ast,\qquad
0\in \partial f(\bar x)+A^\ast \bar y^\ast,\qquad
\langle \bar y^\ast,b-A\bar x\rangle=0.
\]
这就是定义~\ref{def:kkt-cone-system} 的 KKT 系统，故两条等价。
\end{proof}

\begin{corollary}[KKT 的 Fenchel 次微分写法]\label{cor:kkt-via-fenchel-subdifferential}
在定理~\ref{thm:infinite-dimensional-kkt} 的假设下，KKT 系统等价于
\[
-A^\ast \bar y^\ast\in \partial f(\bar x),
\qquad
\bar y^\ast\in \partial \iota_{b-K}(A\bar x).
\]
因此，锥规划中的 KKT 条件本质上并不是另一套平行理论，而是第 \ref{sec:6-4} 节推论~\ref{cor:fenchel-kkt-subdifferential} 在指标函数 $g=\iota_{b-K}$ 上的具体化。
\end{corollary}

\begin{proof}
这正是定理~\ref{thm:infinite-dimensional-kkt} 证明中的中间等价关系。
\end{proof}

\subsection{典型例子}

\begin{example}[状态约束最优控制中的测度乘子]\label{ex:measure-valued-multiplier}
在连续时间最优控制里，若状态约束写成
\[
x(t)\le c(t),\qquad t\in[0,T],
\]
则对应的乘子往往不再是有限维向量，而是区间 $[0,T]$ 上的正测度。把这个例子放在这里，是为了说明定义~\ref{def:kkt-cone-system} 中的 $\bar y^\ast$ 的确可能住在对偶空间甚至测度空间里；“乘子是测度”不是修辞，而是无穷维 KKT 的标准现象。
\end{example}

\begin{example}[线性规划的向量乘子]
当 $K=\mathbb R_+^m$ 时，定义~\ref{def:kkt-cone-system} 退化为
\[
Ax\le b,\qquad y\ge 0,\qquad 0\in \partial f(x)+A^\top y,\qquad y^\top(b-Ax)=0.
\]
把这个例子放在这里，是为了把 Boyd 书中的 KKT 条件完整回收：有限维里熟悉的乘子向量，只是本节抽象对偶变量的一个最简单特例。
\end{example}

\begin{example}[影子价格解释]
在资源配置模型里，约束 $b-Ax\in K$ 表示容量或安全余量非负，而对偶可行变量 $y^\ast\in K^\ast$ 则可解释为影子价格。互补松弛
\[
\langle y^\ast,b-Ax\rangle=0
\]
意味着“有价格的资源必然被用到边界，有剩余的资源不会产生边际价格”。把这个例子放在这里，是为了给 KKT 系统一个管理科学层面的可解释语义。
\end{example}

\subsection*{有限维对照}

Boyd 第 5 章中的 KKT 理论常写成几行矩阵方程，因此读者很容易把它理解成线性代数技巧。本节说明，这些公式真正依赖的是四个抽象结构：锥可行性、对偶锥、次微分驻点条件和互补松弛。第 \ref{sec:8-1} 节会进一步看到，一旦回到 $\mathbb R^n$，很多拓扑和对偶上的麻烦都会自动消失；但正因为本节已经把一般形式写清楚，第 \ref{sec:10-3} 节算法中的原--对偶分裂才能被视作 KKT 系统的计算实现，而不只是某种经验公式。

\subsection*{习题（待补）}

\begin{exercise}[KKT 系统占位]
补一题：对一个给定的锥规划写出定义~\ref{def:kkt-cone-system} 的四条条件，并判断某个候选原--对偶对是否满足它们。
\end{exercise}

\begin{exercise}[乘子空间桥接题占位]
补一题：比较一个有限维线性规划与一个函数空间约束问题中的乘子空间，说明为什么前者的乘子是向量，而后者可能是测度。
\end{exercise}
